Put
\[
 T_*=0.000002349547147305,\qquad \gamma_*=0.8818276493988.
\]
All these displayed terminating decimals denote exact rational numbers.
The certificate checks the following integer comparison:
\begin{equation}\label{eq:tail-rational}
 T_*^2(14\cdot501^7)-C_*^2
 =\frac{3104661541167915364988902768147}
        {20000000000000000000000000000000000}>0.
\end{equation}
Thus (\ref{eq:head-bounds}) and (\ref{eq:tail-final}) give
\begin{align}
 \Gamma
 &> B-h_*-T_*\nonumber\\
 &=0.88182764939885093252
   =\frac{22045691234971273313}{25000000000000000000}
   >\gamma_*,\label{eq:gamma-exact}\\
 B-h_*-T_*-\gamma_*
 &=\frac{1273313}{25000000000000000000}>0.\nonumber
\end{align}
In particular, the coefficient margin is positive with a fixed uniform
constant, independent of the dimensions or entries of $M$.

We give an elementary rational verification of the last reciprocal
comparison, so it does not depend on a decimal approximation to $\pi$.
Machin's identity is
\[
 \pi=16\arctan(1/5)-4\arctan(1/239).
\]
It follows, for example, by applying the tangent addition law to
$4\arctan(1/5)-\arctan(1/239)$ and observing that this angle lies in
$(0,\pi/2)$ and has tangent one. For $0<x<1$, let
$s_{64}(x)=\sum_{k=0}^{63}(-1)^kx^{2k+1}/(2k+1)$. The alternating
series remainder gives
\[
 s_{64}(x)<\arctan x<s_{64}(x)+x^{129}/129.
\]
Substitution at $1/5$ and $1/239$ proves, by rational arithmetic,
\[
 \pi<\pi_U:=3.141592653589793238462643383280.
\]
Finally,
\begin{equation}\label{eq:reciprocal}
 2\gamma_*\,(1.781296297373)-\pi_U
 =\frac{8045255950767709}{12500000000000000000000000000}>0.
\end{equation}
Lemma~\ref{lem:rounding}, (\ref{eq:gamma-exact}), and
(\ref{eq:reciprocal}) prove Theorem~\ref{thm:main}. Taking the supremum
over finite matrices is legitimate because the same constant
$\pi/(2\gamma_*)$ works for every matrix, and that constant is strictly
less than the displayed decimal.

The degree-501 calculation also proves the conservative bound
$K_G^{\R}<1.78130$ using only $C_4<27159$ in place of
(\ref{eq:C4}). Indeed, rational comparisons give
\[
 \frac{27159}{\sqrt{14\cdot501^7}}<0.000002578800327634,
 \quad B-h_*-0.000002578800327634>0.8818274201456,
\]
and $\pi/(2\cdot0.8818274201456)<1.781296760466<1.78130$.
The sharper stated bound uses specifically (\ref{eq:rational-square})
and (\ref{eq:Cmargin}).
